\section*{Capítulo \ref{ch:b2:corrosion} --- Corrosão e proteção}

\begin{solution}{exo:b2:corrosion:1}
$j = \qty{0.020}{A/m^2}$: $de/dt = 0.020 \times 65.38/(2 \times 96\,485 \times 7.13
\times 10^6) = \qty{9.5e-13}{m/s}$, isto é, \qty{0.030}{mm} por ano.
\end{solution}

\begin{solution}{exo:b2:corrosion:2}
Parafusos de aço em cobre: o ferro, o menos nobre, é o ânodo, e os pequenos
parafusos sobre um grande cátodo de cobre corroem depressa. Pregos de zinco no
aço: o zinco corrói e protege o aço em volta dele.
\end{solution}

\begin{solution}{exo:b2:corrosion:3}
Na fresta estreita entre as placas e sob a cabeça do parafuso, onde a água se
renova lentamente e seu oxigênio logo se esgota: por \omterm{def:b2:corrosion:differential}{aeração diferencial},
essas zonas confinadas tornam-se ânodos, e as superfícies abertas, cátodos.
\end{solution}

\begin{solution}{exo:b2:corrosion:4}
A lata estanhada: o estanho, mais nobre que o ferro, faz do aço exposto o
ânodo de um grande cátodo. No balde galvanizado, é o zinco que corrói, no
lugar do aço.
\end{solution}

\begin{solution}{exo:b2:corrosion:5}
A agitação afina a \omterm{def:b2:current-potential-curves:limiting-current}{camada de difusão}: o patamar do oxigênio sobe, a curva
anódica do ferro o encontra mais acima, e assim tanto o potencial quanto a
\omterm{def:b2:corrosion:uniform}{corrente de corrosão} aumentam. Sem oxigênio, a única reação catódica que resta
em água neutra é a lenta redução da água: o \omterm{def:b2:corrosion:uniform}{potencial de corrosão} cai e a
\omterm{def:b2:corrosion:uniform}{corrente de corrosão} se torna muito pequena.
\end{solution}

\begin{solution}{exo:b2:corrosion:6}
$t = 0.90 \times 5000 \times 2 \times 96\,485/(65.38 \times 0.10) = \qty{1.33e8}{s}$,
cerca de 4.2~anos.
\end{solution}

\begin{solution}{exo:b2:corrosion:7}
Carga por quilograma, $nF/M$: zinco \qty{820}{A.h/kg}, magnésio \qty{2.21e3}{A.h/kg},
alumínio \qty{2.98e3}{A.h/kg}. Potenciais padrão: \ce{Zn^2+}/\ce{Zn} $-0.76$,
\ce{Mg^2+}/\ce{Mg} $-2.36$, \ce{Al^3+}/\ce{Al} \qty{-1.68}{V}. Em um solo seco e
resistivo, a corrente precisa atravessar uma grande resistência: o magnésio,
muito mais abaixo do ferro, oferece a maior tensão motriz.
\end{solution}

\begin{solution}{exo:b2:corrosion:8}
$U = 2.0 \times 4.0 + 1.5 = \qty{9.5}{V}$; potência \qty{19}{W}; por ano, $19 \times
8766 = \qty{1.7e5}{W.h}$, \qty{1.7e2}{kW.h}.
\end{solution}

\begin{solution}{exo:b2:corrosion:9}
Na linha d'água, o metal é molhado por água rica em oxigênio, um cátodo; logo
abaixo, a água é menos aerada e a zona se torna o ânodo da pilha formada com
a linha d'água: o aço corrói ali, em uma faixa.
\end{solution}

\begin{solution}{exo:b2:corrosion:10}
Em água aerada, o \omterm{def:b2:corrosion:uniform}{potencial de corrosão} do aço inoxidável cai em seu patamar
passivo: uma corrente minúscula, com o filme se regenerando. Os íons cloreto
rompem o filme em pontos fracos; o metal nu ali é um pequeno ânodo diante de
um enorme cátodo passivo, e sua densidade de corrente é enorme. Dentro do
pite, os íons metálicos se hidrolisam e acidificam a solução, o cloreto é
atraído para equilibrar a carga, e o filme não consegue se refazer: o pite
se aprofunda.
\end{solution}

\begin{solution}{exo:b2:corrosion:11}
O oxigênio é reduzido em todos os \qty{22}{cm^2}: \omterm{def:b2:current-potential-curves:ie-curve}{corrente catódica} $22 \times 5 =
\qty{110}{\micro A}$, toda fornecida pela oxidação dos \qty{2}{cm^2} de aço
perto da junta, $j = \qty{55}{\micro A/cm^2}$. A perda é $5.5$ vezes a do ferro
a \qty{10}{\micro A/cm^2}, cerca de \qty{0.64}{mm} por ano.
\end{solution}

\begin{solution}{exo:b2:corrosion:12}
O ferro é imune abaixo da reta \ce{Fe^2+}/\ce{Fe}, $-0.41 + 0.0296\log 10^{-6} =
\qty{-0.59}{V}$: mantido abaixo de cerca de \qty{-0.6}{V}, o aço deixa de se
dissolver. Muito mais abaixo, a água é reduzida em ritmo crescente sobre o aço
(a reta \ce{H+}/\ce{H2} está em \qty{-0.41}{V} a pH 7, e a \omterm{def:b2:current-potential-curves:overpotential}{sobretensão} sobre o aço é
de alguns décimos de volt): corrente é desperdiçada e o hidrogênio descola o
revestimento.
\end{solution}

\begin{solution}{pb:b2:corrosion:1}
\textbf{1.} \ce{Fe -> Fe^2+ + 2e-}; \ce{O2 + 2H2O + 4e- -> 4OH-}.
\textbf{2.} A redução do oxigênio, lenta e no patamar de difusão, limita a
\omterm{def:b2:current-potential-curves:ie-curve}{corrente catódica}; a curva anódica do ferro a encontra nesse patamar.
\textbf{3.} $j = \qty{0.10}{A/m^2}$: $0.10 \times 3.156 \times 10^7/(2 \times 96\,485)
\times 55.845 = \qty{913}{g}$ por metro quadrado e por ano.
\textbf{4.} $913/7.87 \times 10^6 = \qty{1.16e-4}{m}$, \qty{0.12}{mm} por ano.
\textbf{5.} $4/0.116 = 34$~anos.
\textbf{6.} A corrosão não é uniforme: frestas, diferenças de aeração entre
solos e defeitos do revestimento concentram a \omterm{def:b2:current-potential-curves:ie-curve}{corrente anódica} em pequenas
áreas, onde os pites perfuram a parede muito antes.
\textbf{7.} $E^\circ = \Delta_f G^\circ/(nF)$ para cada par: \ce{Fe} $-0.41$, \ce{Zn} $-0.76$,
\ce{Mg} $-2.36$, \ce{Al} \qty{-1.68}{V}.
\textbf{8.} Zinco, magnésio e alumínio, todos abaixo do ferro.
\textbf{9.} O alumínio se passiva: seu filme de óxido detém sua dissolução, e
ele passa a fornecer pouca corrente, a menos que seja ligado para impedir o
filme (o que funciona na água do mar, não nos solos).
\textbf{10.} A corrente precisa atravessar a resistência do solo: a tensão
motriz entre o ânodo e o aço protegido é de cerca de \qty{1.9}{V} para o
magnésio contra \qty{0.35}{V} para o zinco (a partir dos potenciais padrão).
\textbf{11.} $3.156 \times 10^7 \times 24.305/(2 \times 96\,485 \times 0.50) =
\qty{7.95}{kg}$ por ampère e por ano.
\textbf{12.} $\pi \times 0.30 \times 10\,000 = \qty{9.42e3}{m^2}$.
\textbf{13.} \qty{94.2}{m^2}.
\textbf{14.} $0.020 \times 94.2 = \qty{1.88}{A}$.
\textbf{15.} $1.885 \times 20 \times 3.156 \times 10^7 = \qty{1.19e9}{C}$.
\textbf{16.} $1.19 \times 10^9 \times 24.305/(2 \times 96\,485 \times 0.50) =
\qty{3.00e5}{g}$, \qty{300}{kg} (ou $1.885 \times 20 \times 7.95$).
\textbf{17.} $300/15 = 20$ ânodos.
\textbf{18.} A resistência do solo limita a distância até a qual um único ânodo
consegue levar sua corrente ao longo do tubo: ânodos distribuídos o protegem
de modo uniforme.
\textbf{19.} Um retificador impulsiona corrente de um ânodo inerte enterrado em
um leito anódico, através do solo, até o tubo ligado a seu polo negativo.
\textbf{20.} $1.885 \times 5.0 + 1.5 = \qty{10.9}{V}$.
\textbf{21.} $10.9 \times 1.885 = \qty{20.6}{W}$; $20.6 \times 8766 = \qty{1.8e5}{W.h}$,
\qty{180}{kW.h} por ano.
\textbf{22.} A água é reduzida no tubo: forma-se hidrogênio, que descola o
revestimento e pode fragilizar o aço, e corrente é desperdiçada.
\textbf{23.} $E < -0.41 + 0.0296\log(10^{-6}) = \qty{-0.59}{V}$.
\textbf{24.} Vinte anos de proteção exigem $\boldsymbol{\approx \qty{3.0e2}{kg}}$
de ânodos de magnésio.
\end{solution}
